% Human source transcription, not native author TeX.
% https://ocw.mit.edu/courses/18-238-geometry-and-quantum-field-theory-spring-2023/mit18_238_s23_week02.pdf
% Theorem 2.17（印刷页26）及§2.9；并入§2.1--2.2的Gaussian与插入证明、§2.3的截断设置、Theorem 2.8与Lemma 2.10、Theorem 2.18与Corollary 2.19、Lemma 2.21。对应印刷页15--21、26--28。只计多维驻相问题。

\paragraph{计题范围（助手）.} The main problem is Theorem 2.17. The Gaussian insertion formula, the one-dimensional localization argument, the Morse lemma, and the noncritical estimate below are supporting parts of the same proof, not separate items.

\begin{theorem}[Theorem 2.17: Multidimensional stationary phase formula]
Let $V$ be a real vector space of dimension $d$ with a fixed volume element $dx$, and $D\subset V$ a compact region with smooth boundary. Let $f,g:D\to\R$ be smooth functions. Assume that $f$ has a unique critical point $c$ in $D$, that $c$ is an interior point, that $\det f''(c)\ne0$, and that $g$ has vanishing derivatives of all orders on $\partial D$. Then
\[
\int_D g(x)e^{if(x)/\hbar}\,dx
=\hbar^{d/2}e^{if(c)/\hbar}I(\hbar),\qquad \hbar>0,
\]
where $I$ extends to a smooth function on $[0,\infty)$ such that
\[
I(0)=(2\pi)^{d/2}e^{\pi i\sigma/4}
\frac{g(c)}{\sqrt{|\det f''(c)|}},
\]
and $\sigma$ is the signature of the symmetric bilinear form $f''(c)$.
\end{theorem}

\subsection*{Gaussian integrals: §2.1}
Let $M(V)$ be the set of nondegenerate complex-valued symmetric bilinear forms on $V$ with nonnegative definite real part, and $M^\circ(V)$ the forms with positive definite real part. If $B=P+iQ\in M^\circ(V)$, then $P^{-1}Q$ is self-adjoint with respect to $P$, and diagonalizes in an orthonormal basis. In that basis $B(x,y)=\sum_j a_jx_jy_j$, with $\Re a_j=1$. Thus $B^{-1}\in M^\circ(V^*)$. The map $B\mapsto B^{-1}$ gives a homeomorphism $M(V)\simeq M(V^*)$ restricting to the interiors.

Fix the continuous branch $(\det B)^{-1/2}$ positive on positive definite forms. It exists and is unique, since $M(V)$ is star-like with respect to a point of $M^\circ(V)$. Let $\mathcal S(V)$ be the Schwartz space and $\mathcal S'(V)$ the tempered distributions. The Fourier transform convention is
\[
\mathcal F(g)(p)=(2\pi)^{-d/2}\int_V g(x)e^{-i(p,x)}\,dx.
\]
It extends to tempered distributions. The function $e^{-B(x,x)/2}$ is a tempered distribution for $\Re B\ge0$, depends continuously on $B$, and is Schwartz when $\Re B>0$.

\begin{lemma}[Lemma 2.1: Gaussian integral]
For $B\in M(V)$,
\[
\mathcal F\left(e^{-B(x,x)/2}\right)
=(\det B)^{-1/2}e^{-B^{-1}(p,p)/2}.
\]
\end{lemma}
\begin{proof}
By continuity it suffices to prove this for $\Re B>0$. Diagonalizing reduces to $d=1$, where for $\Re a>0$ we need
\[
\frac1{\sqrt{2\pi}}\int_{-\infty}^{\infty}e^{-ipx-ax^2/2}\,dx
=\frac1{\sqrt a}e^{-p^2/(2a)}.
\]
Both sides are holomorphic in $a$, so it is enough to check real positive $a$. Completing the square gives
\[
\frac{e^{-a^{-1}p^2/2}}{\sqrt{2\pi}}
\int_{-\infty}^{\infty}e^{-a(x+ia^{-1}p)^2/2}\,dx.
\]
By Cauchy's theorem,
\[
\frac1{\sqrt{2\pi}}\int_{\R+ia^{-1}p}e^{-ax^2/2}\,dx
=\frac1{\sqrt{2\pi}}\int_\R e^{-ax^2/2}\,dx.
\]
The result follows by rescaling the Poisson integral $\int_\R e^{-x^2}\,dx=\sqrt\pi$.
\end{proof}
In the sense of the Gaussian formula,
\[
(2\pi)^{-d/2}\int_V e^{-B(x,x)/2}\,dx=(\det B)^{-1/2}.
\]

\subsection*{Gaussian integrals with insertions: §2.2}
For $g\in\mathcal S(V)$, set
\[
I_g(\hbar)=\int_V g(\hbar^{1/2}x)e^{-B(x,x)/2}\,dx\quad(\hbar>0),
\qquad I_g(0)=(2\pi)^{d/2}(\det B)^{-1/2}g(0).
\]
For a basis $e_1,\ldots,e_d$ and its dual, let $\Delta_B=\sum_{j=1}^d\partial_{B^{-1}e_j^*}\partial_{e_j}$.

\begin{theorem}[Theorem 2.2]
For $\hbar\ge0$,
\[
I'_g(\hbar)=I_{\frac12\Delta_Bg}(\hbar).
\]
Thus $I_g\in C^\infty[0,\infty)$. In particular, if $g$ vanishes at the origin to order $2n+1$, then $I_g(0)=\cdots=I_g^{(n)}(0)=0$.
\end{theorem}

\begin{lemma}[Lemma 2.3]
$I_g$ is continuous.
\end{lemma}
\begin{proof}
Only continuity at zero requires proof. By Plancherel's theorem and Lemma 2.1,
\[
\begin{aligned}
I_g(\hbar)
&=\bigl(g(\hbar^{1/2}x),e^{-B(x,x)/2}\bigr)\\
&=\hbar^{-d/2}(\det B)^{-1/2}
\bigl(\widehat g(\hbar^{-1/2}p),e^{-B^{-1}(p,p)/2}\bigr)\\
&=(\det B)^{-1/2}\bigl(\widehat g(p),e^{-\hbar B^{-1}(p,p)/2}\bigr).
\end{aligned}
\]
But $e^{-\hbar B^{-1}(p,p)/2}\to1$ in $\mathcal S'(V^*)$ as $\hbar\to0$, by continuous dependence of the complex Gaussian distribution on the form. Therefore
\[
\lim_{\hbar\to0}I_g(\hbar)
=(\det B)^{-1/2}(\widehat g,1)
=(2\pi)^{d/2}(\det B)^{-1/2}g(0)=I_g(0).
\]
\end{proof}

\begin{lemma}[Lemma 2.4]
If $\ell\in V^*$ and $f\in\mathcal S(V)$, then
\[
I_{\ell f}(\hbar)=\hbar I_{\partial_{B^{-1}\ell}f}(\hbar).
\]
\end{lemma}
\begin{proof}
We have
\[
\begin{aligned}
I_{\ell f}(\hbar)
&=\hbar^{1/2}\bigl(\ell(x)f(\hbar^{1/2}x),e^{-B(x,x)/2}\bigr)\\
&=\hbar^{1/2}\bigl(f(\hbar^{1/2}x),\ell(x)e^{-B(x,x)/2}\bigr)\\
&=-\hbar^{1/2}\bigl(f(\hbar^{1/2}x),\partial_{B^{-1}\ell}e^{-B(x,x)/2}\bigr)\\
&=\hbar^{1/2}\bigl(\partial_{B^{-1}\ell}[f(\hbar^{1/2}x)],e^{-B(x,x)/2}\bigr)\\
&=\hbar\bigl((\partial_{B^{-1}\ell}f)(\hbar^{1/2}x),e^{-B(x,x)/2}\bigr)
=\hbar I_{\partial_{B^{-1}\ell}f}(\hbar).
\end{aligned}
\]
\end{proof}

\begin{proof}[Proof of Theorem 2.2]
For $\hbar>0$, direct differentiation gives
\[
I'_g(\hbar)=\tfrac12\hbar^{-1}I_{Eg}(\hbar),\qquad
E=\sum_{j=1}^d e_j^*\partial_{e_j}.
\]
By Lemma 2.4 this is $I_{\frac12\Delta_Bg}(\hbar)$. By Lemma 2.3 it suffices to prove $I_g\in C^1[0,\infty)$; repeated application gives smoothness.
For a positive definite form $C$,
\[
I_{e^{-C(x,x)/2}}(\hbar)
=\int_V e^{-(B+\hbar C)(x,x)/2}\,dx
=(2\pi)^{d/2}\det(B+\hbar C)^{-1/2},
\]
which is analytic, hence continuously differentiable on $[0,\infty)$. Subtracting a multiple of this Gaussian from $g$, it suffices to consider $g(0)=0$.
Such $g$ is a linear combination of functions $\ell f$, where $f\in\mathcal S(V)$ and $\ell\in V^*$. For $g=\ell f$, Lemma 2.4 gives
\[
I'_g(0)=I_{\partial_{B^{-1}\ell}f}(0)=I_{\frac12\Delta_Bg}(0),
\]
since
\[
\tfrac12\Delta_B(\ell f)(0)
=\sum_j\ell(e_j)\partial_{B^{-1}e_j^*}f(0)
=\partial_{B^{-1}\ell}f(0).
\]
This completes the proof.
\end{proof}

\subsection*{Localization and the one-dimensional argument: §2.3--§2.4}
\paragraph{截取说明（助手）.} The cutoff construction in the next paragraph is the part of the proof of Theorem 2.6 expressly reused by Theorem 2.8. The subsequent lemma and proof are retained together; no steepest-descent problem is separately counted.

Take coordinates with $c=0$, $f(c)=0$, and $f(x)=Mx^2/2$ on a neighborhood $U$ of zero, where $M=f''(c)$. Let $h$ be a smooth bump supported in $U$ and equal to one in a smaller neighborhood of zero. Write $g=g_1+g_2$ with $g_1=hg$, $g_2=(1-h)g$. The normalized integral splits accordingly as $I=I_1+I_2$. For the local part extend $g_1$ by zero to the real line and use $y=\hbar^{-1/2}x$.

\begin{lemma}[Lemma 2.10: Riemann lemma]
Let $f:[a,b]\to\R$ be smooth, with $f'(x)>0$ throughout. If $g\in C^n[a,b]$ satisfies
\[
g(a)=\cdots=g^{(n-1)}(a)=g(b)=\cdots=g^{(n-1)}(b)=0,
\]
then $J(\hbar):=\int_a^b g(x)e^{if(x)/\hbar}\,dx=O(\hbar^n)$ as $\hbar\to0$.
If $g$ is smooth and all endpoint derivatives vanish, then $J$, extended by $J(0)=0$, is smooth on $[0,\infty)$ and all its derivatives are rapidly decaying as $\hbar\to0$.
\end{lemma}
\begin{proof}
Changing variables reduces to $f(x)=x$. Induct on $n$; the case $n=0$ is obvious. Integration by parts gives
\[
\int_a^b g(x)e^{ix/\hbar}\,dx
=i\hbar\int_a^b g'(x)e^{ix/\hbar}\,dx,
\]
which proves the induction step. The second assertion follows by repeated differentiation and the first assertion.
\end{proof}

\begin{theorem}[Theorem 2.8: one-dimensional supporting case]
If $f,g$ are smooth on $[a,b]$, $f$ has a unique critical point $a<c<b$ with $f''(c)\ne0$, and all derivatives of $g$ vanish at the endpoints, then
\[
\int_a^b g(x)e^{if(x)/\hbar}\,dx
=\hbar^{1/2}e^{if(c)/\hbar}I(\hbar),
\]
where $I\in C^\infty[0,\infty)$ and
\[
I(0)=\sqrt{2\pi}\,e^{\pm\pi i/4}\frac{g(c)}{\sqrt{|f''(c)|}}.
\]
The sign is the sign of $f''(c)$.
\end{theorem}
\begin{proof}
As in the steepest-descent argument, reduce to $c=0$, $f=Mx^2/2$ near zero, and split with the cutoff above. By Lemma 2.10(ii), the contribution of $g_2$ is rapidly decaying together with all derivatives, so it suffices to prove the result for $g=g_1$. Following the same rescaling,
\[
I(\hbar)=\int_{-\infty}^{\infty}g(\hbar^{1/2}y)e^{iMy^2/2}\,dy.
\]
The assertion follows from the Gaussian value at zero and Theorem 2.2.
\end{proof}

\subsection*{Morse lemma: §2.8}
\begin{theorem}[Theorem 2.18: Separation of variables]
Let $f$ be smooth on an open ball $0\in B\subset\R^d$, with a nondegenerate critical point at zero and $f(0)=0$. There are local coordinates near zero in which
\[
f(x_1,\ldots,x_d)=f(x_1,\ldots,x_{d-1},0)\pm x_d^2.
\]
\end{theorem}
\paragraph{坐标排印说明（助手）.} In the following source proof, $u=g(y,v)$ is the inverse coordinate change. Accordingly the displayed substitution is written $f_*(y,v)=f(y,g(y,v))$; the printed right side has $g(y,u)$. The chain-rule line is written with its indicated evaluation, and the order of $(y,u)$ is kept consistent. These are notation corrections, not an alternative proof.
\begin{proof}
By a linear change of variables, the quadratic part can be written $Q(y)\pm u^2$, where $y=(x_1,\ldots,x_{d-1})$, $u=x_d$. Consider the hypersurface $S$ given by $\partial_u f(y,u)=0$.
The linear part of $\partial_u f(y,u)$ is $\pm2u$. By the implicit function theorem there is a change of coordinates $F$ near zero with $dF(0)=1$, replacing $u$ by $v=\pm\tfrac12\partial_u f(y,u)$ and keeping $y$ unchanged. Thus $u=g(y,v)$ with $(\partial_vg)(0,0)\ne0$. Put $f_*(y,v)=f(y,g(y,v))$. By the chain rule,
\[
\partial_v f_*(y,v)
=(\partial_u f)(y,u)\frac{\partial u}{\partial v}
=\pm2v\partial_vg(y,v).
\]
Thus $S$ is given in the new coordinates by $v=0$. We may therefore assume from the start that $S$ is given by $u=0$. Then $(\partial_u f)(y,0)=0$, and
\[
f(y,u)-f(y,0)=h(y,u)u^2,
\]
with $h$ smooth and $h(0,0)=\pm1$. Replacing $u$ by $\widetilde u=\sqrt{|h(y,u)|}\,u$ and keeping $y$ unchanged makes $h=\pm1$. Hence
\[
f(y,u)=f(y,0)\pm u^2,
\]
as claimed.
\end{proof}

\begin{corollary}[Corollary 2.19: Morse lemma]
Under the same hypotheses, local coordinates can be chosen with
\[
f=x_1^2+\cdots+x_m^2-x_{m+1}^2-\cdots-x_d^2.
\]
In other words, near a nondegenerate critical point a smooth function is equivalent by a coordinate change to its quadratic part.
\end{corollary}
\begin{proof}
This follows from Theorem 2.18 by induction on dimension.
\end{proof}

\subsection*{Proof of the multidimensional formula: §2.9}
The proofs of the multidimensional steepest-descent and stationary-phase formulas are parallel to their one-dimensional versions, using the Morse lemma. Namely, the Morse lemma permits us to assume that $f$ is quadratic near the critical point. After this, the proof of the steepest-descent formula is identical to the one-variable case. The same applies to the stationary-phase formula, using the following multivariable analog of the Riemann lemma.

\begin{lemma}[Lemma 2.21]
Let $f,g:D\to\R$ be smooth, with all derivatives of $g$ zero on $\partial D$, and with $df$ nowhere zero on $\operatorname{supp}g$. Then
\[
J(\hbar):=\int_D g(x)e^{if(x)/\hbar}\,dx
\]
extends to a smooth function on $[0,\infty)$ and has rapidly decaying derivatives of all orders as $\hbar\to0$.
\end{lemma}
\begin{proof}
Since $df$ is nonzero on $\operatorname{supp}g$, cover the support by local charts $U_i$ in which $f(x)$ is the last coordinate $x_d$. By compactness the cover is finite. Use a partition of unity $h_i$ subordinate to the cover and replace $g$ by $h_i g$, reducing to support in one chart. Change variables so that $f(x)=x_d$. Integrating out $x_1,\ldots,x_{d-1}$ reduces to the one-dimensional case of Lemma 2.10.
\end{proof}

\paragraph{关于原文证明的承接（助手）.}
§2.9以“与一维证明平行”的方式结束主证明；上文已并入它实际指向的截断、重缩放、Gaussian插入公式和Morse坐标构造。这里保留这种引用式组织，不将其改写成另一份独立证明。主公式的系数取自Theorem 2.17原式；平方根分支按§2.1在实正定形式上为正的定义。
